\chapter{Code Generation: From Python to C++ Sweeps}\label{r02:ch:codegen}

waLBerla's fast kernels are not part of its source tree. They are generated during the
build from Python descriptions. This chapter follows that path from the generator
script to the compiled sweep, and points out where the hardware target is chosen.
\Cref{r02:fig:build-flow} is the map for this chapter and for the rest of the report.

\begin{figure}[htbp]
  \centering
  \resizebox{\linewidth}{!}{% From CMake configuration to an object file: where the target and the ISA are decided.
% Orange = changed for RVV, teal = new for RVV, white = unchanged. Markers 1-5 are the
% decision points of Chapter 6.
\begin{tikzpicture}[x=1cm,y=1cm,
    fb/.style={r02box,text width=2.95cm},
    lbl/.style={r02label,minimum height=0pt}]
  % ---- configure time
  \node[fb,text width=3.25cm] (cache) at (-0.15,0)
    {\textbf{CMake cache}\\\texttt{\scriptsize OPTIMIZE\_FOR\_LOCALHOST}\\
     \texttt{\scriptsize DOUBLE\_ACCURACY}, \dots};
  \node[fb,r02changed] (copts) at (4.1,0.75)
    {\textbf{compileroptions/}\\\texttt{GNU.cmake}, \texttt{Clang.cmake}};
  \node[fb,text width=3.3cm] (flags) at (8.2,0.75)
    {\texttt{CMAKE\_CXX\_FLAGS}\\\texttt{\scriptsize -march=rv64gcv\_zvl256b}};
  \node[fb,text width=3.3cm] (cfg) at (4.1,-0.85)
    {\texttt{SweepgenConfig.py}\\{\scriptsize\cls{WalberlaBuildConfig}}};
  \draw[r02arrow] (cache.east) -- ++(0.35,0) |- (copts.west);
  \draw[r02arrow] (cache.east) -- ++(0.35,0) |- (cfg.west);
  \draw[r02arrow] (copts) -- (flags);

  \draw[psgray!60,dashed] (-1.6,-1.85) -- (13.95,-1.85);
  \node[lbl,anchor=south west] at (-1.6,-1.83) {configure time};
  \node[lbl,anchor=north west] at (-1.6,-1.87) {build time};

  % ---- build time: code generation
  \node[fb,r02new] (script) at (0,-3.15)
    {\textbf{generator script}\\\texttt{SweepgenLbmCPU.py}\\(\texttt{gen\_config})};
  \node[fb] (sweepgen) at (4.1,-3.15)
    {\textbf{sweepgen}\\\cls{Sweep}, \cls{LbmBulk}\\target $=$ \code{CurrentCPU}};
  \node[fb,r02changed] (ps) at (8.2,-3.15)
    {\textbf{\ps{}}\\\code{CurrentCPU} $\rightarrow$ \code{RISCV_RVV}\\\cls{RvvCpu} platform};
  \node[fb] (sfg) at (12.3,-3.15)
    {\textbf{pystencils-sfg}\\writes \texttt{gen/*.hpp}, \texttt{gen/*.cpp}};
  \draw[r02arrow] (script) -- (sweepgen);
  \draw[r02arrow] (sweepgen) -- (ps);
  \draw[r02arrow] (ps) -- (sfg);
  \draw[r02arrow] (cfg.south) -- (sweepgen.north);

  % ---- build time: compilation
  \node[r02code,text width=2.95cm] (src) at (12.3,-5.2)
    {\#include ".../rvv.hpp"\\\_\_riscv\_vle64\_v\_f64m1\\\_\_riscv\_vfadd\_vv\_...};
  \node[fb] (cxx) at (8.2,-5.2) {\textbf{C++ compiler}\\GCC 14 / Clang 18};
  \node[r02code,text width=2.95cm] (obj) at (4.1,-5.2) {vle64.v\quad vfadd.vv\\vfmul.vv\quad vse64.v};
  \node[fb] (app) at (0,-5.2) {\textbf{application}\\links \code{walberla::v8}};
  \draw[r02arrow] (sfg) -- (src);
  \draw[r02arrow] (src) -- (cxx);
  \draw[r02arrow] (cxx) -- (obj);
  \draw[r02arrow] (obj) -- (app);
  \draw[r02arrow] (flags.east) -- ++(0.55,0) |- ([yshift=4mm]cxx.east)
    node[lbl,pos=0.15,right,align=left] {global\\flags};

  % ---- decision markers
  \node[r02marker] at ([xshift=1.5mm,yshift=1.2mm]copts.north east) {1};
  \node[r02marker] at ([xshift=1.5mm,yshift=1.2mm]sweepgen.north east) {2};
  \node[r02marker] at ([xshift=1.5mm,yshift=1.2mm]ps.north east) {3};
  \node[r02marker] at ([xshift=1.5mm,yshift=1.2mm]src.north east) {4};
  \node[r02marker] at ([xshift=1.5mm,yshift=1.2mm]script.north east) {5};
\end{tikzpicture}
}
  \caption[From CMake configuration to RVV object code]{From the CMake configuration to
    RVV object code. Configure time (top) produces the global compiler flags and the
    project configuration module; build time (bottom) runs the generator script through
    sweepgen, \ps{} and pystencils-sfg, and compiles the result. Orange boxes were
    changed for RVV, the teal box is new, white boxes are unchanged. The numbered
    markers are the five decision points of \cref{r02:sec:gap-chain}.}
  \label{r02:fig:build-flow}
\end{figure}

\section{The tool chain}\label{r02:sec:toolchain}

Four Python packages from the \emph{pycodegen} family work together:
\begin{description}[style=nextline]
  \item[\ps{}~\cite{pystencils}] compiles a symbolic stencil description (SymPy
    assignments on abstract fields) into one C/C++ function for a given
    \term{target}: generic C, x86 SSE/AVX/AVX-512, ARM NEON/SVE, CUDA, HIP, SYCL, and
    now RISC-V RVV. Report~1 of this collection describes it in detail.
  \item[lbmpy~\cite{lbmpy}] derives lattice Boltzmann methods symbolically (stencil,
    collision operator, equilibrium, forcing, boundary conditions) and hands the
    resulting assignment collections to \ps{}.
  \item[pystencils-sfg~\cite{pystencils-sfg}] (\emph{source file generator}) writes
    complete C++ header and source files around \ps{} kernels: namespaces, classes,
    methods, and the glue that extracts data pointers, shapes and strides from C++
    data structures and passes them to the kernel. It also provides the CMake function
    \code{pystencilssfg_generate_target_sources} that runs generator scripts at build
    time.
  \item[sweepgen] is waLBerla's layer on top of pystencils-sfg
    (\file{sweepgen/src/sweepgen/}). It knows waLBerla's data structures: blocks, the
    classic \cls{GhostLayerField} and the v8 \code{memory::Field}, cell intervals,
    block geometry. It provides \cls{Sweep} (one kernel as a sweep class),
    \cls{LbmBulk} (stream-collide and initialization sweeps for an lbmpy method),
    boundary-condition generators, pack infos for communication, and the build
    configuration object \cls{WalberlaBuildConfig}.
\end{description}

\section{Two code generation paths}\label{r02:sec:two-paths}

waLBerla v7.2 contains two independent code generation layers:
\begin{description}[style=nextline]
  \item[Legacy: \file{python/pystencils_walberla} and \file{lbmpy_walberla}] Enabled with
    \cmakeopt{WALBERLA_BUILD_WITH_CODEGEN}, used by the older applications in
    \file{apps/} and by the classic \file{lbm_generated} module. It is written against
    the \ps{} \emph{1.x} API: its helper \code{get_vectorize_instruction_set(ctx)}
    imports \code{pystencils.backends.simd_instruction_sets}, a module that no longer
    exists in \ps{}~2.0, and picks the newest instruction set the 1.x code knows
    (\code{'sse'}, \code{'avx'}, \code{'avx512'}, \code{'neon'}, \dots). It cannot work
    with the 2.0-based RVV backend without a port.
  \item[Current: sweepgen] Enabled with \cmakeopt{WALBERLA_ENABLE_SWEEPGEN}, built on
    \ps{}~2.0 and pystencils-sfg, used by the v8 examples and tests. This is where new
    development happens, and the only path this work targets.
\end{description}

\section{A generator script}\label{r02:sec:sweepgen-script}

A generator script is an ordinary Python program that opens a
\code{SourceFileGenerator} context and emits code into it.
\Cref{r02:lst:genscript} is the core of \file{SweepgenLbmCPU.py}, the generator of the
benchmark application. The lines that request explicit vectorization are the only
RVV-related part of it; everything else is how any waLBerla LBM application looks.

\rneedspace{12\baselineskip}
\begin{pycode}[nolol=false,caption={[Generator script of the benchmark]Generator script of the benchmark (\texttt{apps/benchmarks/SweepgenLbmCPU/SweepgenLbmCPU.py}, abridged).},label=r02:lst:genscript]
with SourceFileGenerator(keep_unknown_argv=True) as sfg:
    args = parser.parse_args(sfg.context.argv)      # --variant, --cse-global

    Sweep.use_v8core_fields()                       # generate for walberla::v8::memory::Field
    sfg.namespace("SweepgenLbmCPU::gen")

    if args.variant == "explicit":
        lbm_gen_config = ps.CreateKernelConfig(target=ps.Target[args.explicit_target])  # CurrentCPU
        lbm_gen_config.cpu.vectorize.enable = True
        #   waLBerla fields are fzyx with x fastest: the innermost loop has stride 1
        lbm_gen_config.cpu.vectorize.assume_inner_stride_one = True
    else:
        lbm_gen_config = ps.CreateKernelConfig(target=ps.Target.GenericCPU)

    lbm_config = LBMConfig(stencil=LBStencil(Stencil.D3Q19), method=Method.SRT,
                           compressible=True, relaxation_rate=sp.Symbol("omega"))
    lbm_optimisation = LBMOptimisation(simplification="auto", cse_global=args.cse_global,
                                       field_layout="fzyx")

    lbm_bulk = LbmBulk(sfg, "LBM", lbm_config, lbm_optimisation, lbm_gen_config)
    sfg.generate(lbm_bulk)                          # -> classes LBM::StreamCollide, LBM::InitPdfs

    #   Initial state (tanh, Piecewise): not vectorizable, never timed -> always generic
    sfg.generate(Sweep("SetInitialState", initial_state_assignments,
                       ps.CreateKernelConfig(target=ps.Target.GenericCPU)))
\end{pycode}

The script runs in the build, once per executable, with the arguments given in
\file{CMakeLists.txt} (\cref{r02:sec:app}). \code{sfg.context.argv} carries those
arguments; pystencils-sfg consumes its own options (output directory, configuration
module) first.

\section{The build configuration and the target}\label{r02:sec:buildconfig}

Every \cls{Sweep} and \cls{LbmBulk} starts from the code generator configuration that
the build implies, and then applies the script's own configuration on top
(\code{gen_config.override(...)}). The build configuration is the
\cls{WalberlaBuildConfig} object returned by \code{project_info()} in
\file{SweepgenConfig.py} (\cref{r02:lst:sweepgen-config}). Its two key methods are in
\cref{r02:lst:buildconfig}.

\rneedspace{12\baselineskip}
\begin{pycode}[nolol=false,caption={[Target selection in sweepgen's build configuration]Target selection in \texttt{sweepgen/src/sweepgen/build\_config.py} (unchanged upstream code, abridged).},label=r02:lst:buildconfig]
@property
def target(self) -> Target:
    if self.override.target is not None:
        target = self.override.target
    elif self.cuda_enabled:
        target = Target.CUDA
    elif self.hip_enabled:
        target = Target.HIP
    else:
        #  CPU target
        if self.optimize_for_localhost:
            target = Target.CurrentCPU      # "the best this host can do"
        else:
            target = Target.GenericCPU      # portable C++
    ...
    return target

def get_pystencils_config(self) -> CreateKernelConfig:
    cfg = CreateKernelConfig()
    cfg.default_dtype = Fp(64) if self.use_double_precision else Fp(32)
    cfg.jit = no_jit                        # generate only, never compile in Python
    cfg.target = self.target
    if self.openmp_enabled:
        cfg.cpu.openmp.enable = True
    return cfg
\end{pycode}

Two facts follow from this code:
\begin{enumerate}
  \item On a CPU build, the target is \code{Target.CurrentCPU} exactly when
        \cmakeopt{WALBERLA_OPTIMIZE_FOR_LOCALHOST} is on. \ps{} resolves
        \code{CurrentCPU} when the kernel is created, by asking the host:
        \code{CreateKernelConfig.get_target()} calls \code{Target.auto_cpu()}, which
        returns the last (most capable) entry of \code{_available_vector_targets()}, or
        \code{GenericCPU} if the list is empty. On an AVX2 workstation that is
        \code{X86_AVX}; on the K1 it is now \code{RISCV_RVV}
        (\cref{r02:sec:ps-target}). Code is generated on the machine that builds, so
        the build host must be the target machine, as with \code{-march=native}.
  \item \code{get_pystencils_config} never sets \code{cpu.vectorize.enable}. A vector
        target without that switch only changes the platform object inside \ps{}; the
        loops stay scalar. This is the same on x86: waLBerla sweeps are scalar unless a
        generator script requests vectorization, as \cref{r02:lst:genscript} does.
\end{enumerate}

\section{Anatomy of a generated sweep}\label{r02:sec:generated-sweep}

For each \cls{Sweep}, sweepgen generates a C++ class in the header and its methods in
the source file. The header of \code{LBM::StreamCollide} (generated on the K1, double
precision) contains:
\begin{itemize}
  \item a \code{detail::StreamCollideProperties} struct that holds the fields
        (\code{f_}, \code{rho_}, \code{u_}) and scalar parameters (\code{omega_});
  \item a constructor taking those, getters for the fields, and getter/setter pairs for
        the parameters, so an application can change \code{omega} between steps;
  \item \code{operator()(IBlock*)} for a whole block, \code{runOnCellInterval(...)} for
        a sub-region (used for communication hiding), and \code{swapShadowBuffers};
  \item a \code{ShadowBufferCache} member that owns the temporary PDF field
        \code{f_tmp} per block, so the application never sees it.
\end{itemize}
The source file contains the \ps{} kernel as a free function in a namespace
\code{StreamCollide_kernels}, followed by the method bodies.
\Cref{r02:lst:generated-wrapper} shows how \code{operator()} maps the v8 field views to
the kernel's raw parameters.

\rneedspace{12\baselineskip}
\lstinputlisting[style=cpp,nolol=false,basicstyle=\ttfamily\scriptsize,
  caption={[A generated sweep's operator()]Generated \texttt{StreamCollide::operator()} (explicit variant, double precision; abridged, comments added in \texttt{//}).},
  label=r02:lst:generated-wrapper]{streamcollide-wrapper.cpp}

The kernel signature differs between variants in one revealing detail. The generic
(scalar and auto-vectorized) kernel receives the stride along $x$ as a run-time argument
\code{_stride_f_0}, so the compiler cannot know that consecutive cells are adjacent in
memory. With \code{assume_inner_stride_one}, \ps{} fixes that stride to 1 at code
generation time; the parameter disappears, and the wrapper only records the assumption
in a comment. Nothing checks it at run time. For waLBerla fields it holds by
construction (\code{fzyx}, $x$ fastest), but it is a contract, not a guarantee
(\cref{r02:sec:lim-stride}).
